In this section, we provide additional details on the experimental setup used in
the main paper, including the training details of the models,
and the datasets used.

\paragraph{ResNet-9 on CIFAR-10.} We use the ResNet-9 architecture from
\citep{jordan202494}, with the hyperparameters given in
Table~\ref{tab:resnet9_cifar10_hyperparams}. To give concrete details: the
training set $S$ comprises $50,000$ CIFAR-10 training samples, the
learning algorithm
$\mathcal{A}$ is standard supervised training, and we consider 50 measurement
functions $\phi_i$ corresponding to loss on 50 different CIFAR-10 test samples.

\begin{table}[h]
  \centering
  \begin{tabular}{ll}
    \toprule
    Hyperparameter & Value \\
    \midrule
    Learning rate & 1.2 \\
    Weight decay & 0.001 \\
    Bias scale & 8.0 \\
    Batch size & 1000 \\
    Epochs & 12 \\
    Final layer scale & 0.04 \\
    Momentum & 0.875 \\
    Pooling type & Log-sum-exp \\
    Pooling $\varepsilon$ & 0.1 \\
    Width multiplier & 2.5 \\
    LR schedule & One-cycle Linear \\
    LR start multiplier & 0.07 \\
    LR end multiplier & 0.2 \\
    LR peak time & 0.5 \\
    \bottomrule
  \end{tabular}
  \caption{Hyperparameters for ResNet-9 on CIFAR-10.}
  \label{tab:resnet9_cifar10_hyperparams}
\end{table}

\paragraph{Gemma-2B LoRA on IFT Data.} We use the variant of
LESS~\citep{xia2024less} from \citet{engstrom2025optimizing}. In particular, the
training dataset consists of the four instruction fine-tuning sets seen in
Table~\ref{tab:ift_training_dataset} as in LESS. The total number of points is
around $300,000$ and is exactly four combined IFT datasets (Flan
  V2~\citep{longpre2023flan}, CoT~\citep{wei2022chain},
  DOLLY~\citep{conover2023free}, and Open Assistant 1
\citep{kopf2024openassistant}). We test on a randomly chosen (task balanced)
subset of of MMLU comprising 32 test samples. We use 4-shot in-context learning
for these samples. We adapt a LoRA to a Gemma-2B model (the pretraining-only
Gemma-2B model~\citep{team2024gemma}) using the LoRA configuration from
\citet{xia2024less}. For model training, we use the same setup as
\citet{engstrom2025optimizing}, except with $\epsilon_\textrm{root}=10^{-6}$. In
particular, we train with ADAM ($\beta_1=0.95, \beta_2=0.975$, decoupled weight
decay as $10^{-5}$) and a one-cycle linear schedule starting at $10^{-6}$ of the
maximum learning rate, reaching the peak over 25\% of training, then ending at
$0.1$ of the maximum learning rate ($0.0004$). We insert a positive
$\epsilon_\textrm{root}$ into the inverse square root term in the ADAM update to
prevent metagradient (and to a lesser extent update) blowup.

\paragraph{GPT-2 fine-tuning on Wikitext.} We optimize a pre-trained
GPT2~\citep{radford2019language} model on Wikitext~\citep{foundation2022english}
using causal language modeling. We split the Wikitext samples into size 512
context length chunks and into train and test splits, with 256 samples in the
test split and 4608 samples in the train split. We attribute on the test split,
and use 4 epochs of the train split during training. We use the same ADAM
optimizer setup above except that we set $\epsilon_\textrm{root}=10^{-8}$, max
learning rate to $0.0008$, and do not anneal ADAM $\epsilon_{\mathrm{root}}$.

\begin{table}[h]
  \small
  \centering
  \caption{Details of IFT training datasets.}
  \label{tab:ift_training_dataset}
  \begin{tabular}{l l l c c}
    \toprule
    Dataset & \# Instance & Sourced from & Prompt Len. & Complet. Len. \\
    \midrule
    FLAN V2 & 100,000 & NLP datasets and human-written instructions &
    355.7 & 31.2 \\
    CoT & 100,000 & NLP datasets and human-written CoTs & 266 & 53.2 \\
    Dolly & 15,011 & Human-written from scratch & 118.1 & 91.3 \\
    Open Ast. 1 & 55,668 & Human-written from scratch & 34.8 & 212.5 \\
    \bottomrule
  \end{tabular}
\end{table}

